\section{Jacobson schemes}
\label{section-jacobson}

\noindent
Recall that a space is said to be {\it Jacobson} if the closed points are
dense in every closed subset, see
Topology, Section \ref{topology-section-space-jacobson}.

\begin{definition}
\label{definition-jacobson}
A scheme $S$ is said to be {\it Jacobson} if its underlying topological
space is Jacobson.
\end{definition}

\noindent
Recall that a ring $R$ is Jacobson if every radical ideal of $R$
is the intersection of maximal ideals, see
Algebra, Definition \ref{algebra-definition-ring-jacobson}.

\begin{lemma}
\label{lemma-affine-jacobson}
An affine scheme $\Spec(R)$ is Jacobson if and only if
the ring $R$ is Jacobson.
\end{lemma}

\begin{proof}
This is Algebra, Lemma \ref{algebra-lemma-jacobson}.
\end{proof}

\noindent
Here is the standard result characterizing Jacobson schemes.
Intuitively it claims that Jacobson $\Leftrightarrow$ locally Jacobson.

\begin{lemma}
\label{lemma-locally-jacobson}
Let $X$ be a scheme. The following are equivalent:
\begin{enumerate}
\item The scheme $X$ is Jacobson.
\item The scheme $X$ is ``locally Jacobson'' in the sense of
Definition \ref{definition-locally-P}.
\item For every affine open $U \subset X$ the ring $\mathcal{O}_X(U)$
is Jacobson.
\item There exists an affine open covering $X = \bigcup U_i$ such that
each $\mathcal{O}_X(U_i)$ is Jacobson.
\item There exists an open covering $X = \bigcup X_j$
such that each open subscheme $X_j$ is Jacobson.
\end{enumerate}
Moreover, if $X$ is Jacobson then every open subscheme
is Jacobson.
\end{lemma}

\begin{proof}
The final assertion of the lemma holds by
Topology, Lemma \ref{topology-lemma-jacobson-inherited}.
The equivalence of (5) and (1) is
Topology, Lemma \ref{topology-lemma-jacobson-local}.
Hence, using Lemma \ref{lemma-affine-jacobson},
we see that (1) $\Leftrightarrow$ (2).
To finish proving the lemma it suffices to show that ``Jacobson''
is a local property of rings, see Lemma \ref{lemma-locally-P}.
Any localization of a Jacobson ring at an element is Jacobson, see
Algebra, Lemma \ref{algebra-lemma-Jacobson-invert-element}.
Suppose $R$ is a ring, $f_1, \ldots, f_n \in R$ generate the unit
ideal and each $R_{f_i}$ is Jacobson. Then we see that
$\Spec(R) = \bigcup D(f_i)$ is a union of open subsets
which are all Jacobson, and hence $\Spec(R)$ is Jacobson
by Topology, Lemma \ref{topology-lemma-jacobson-local} again.
This proves the second property of Definition \ref{definition-property-local}.
\end{proof}

\noindent
Many schemes used commonly in algebraic geometry are Jacobson, see
Morphisms, Lemma \ref{morphisms-lemma-ubiquity-Jacobson-schemes}.
We mention here the following interesting case.

\begin{lemma}
\label{lemma-complement-closed-point-Jacobson}
Examples of Noetherian Jacobson schemes.
\begin{enumerate}
\item If $(R, \mathfrak m)$ is a Noetherian local ring, then
the punctured spectrum $\Spec(R) \setminus \{\mathfrak m\}$
is a Jacobson scheme.
\item If $R$ is a Noetherian ring with Jacobson radical $\text{rad}(R)$
then $\Spec(R) \setminus V(\text{rad}(R))$ is a Jacobson scheme.
\item If $(R, I)$ is a Zariski pair (More on Algebra, Definition
\ref{more-algebra-definition-zariski-pair})
with $R$ Noetherian, then $\Spec(R) \setminus V(I)$ is a
Jacobson scheme.
\end{enumerate}
\end{lemma}

\begin{proof}
Proof of (3). Observe that $\Spec(R) - V(I)$ has a covering by the affine
opens $\Spec(R_f)$ for $f \in I$. The rings $R_f$ are Jacobson by
More on Algebra, Lemma
\ref{more-algebra-lemma-noetherian-zariski-jacobson-complement}.
Hence $\Spec(R) \setminus V(I)$ is Jacobson by
Lemma \ref{lemma-locally-jacobson}.
Parts (1) and (2) are special cases of (3).

\medskip\noindent
Direct proof of case (1).
Since $\Spec(R)$ is a Noetherian scheme,
$S$ is a Noetherian scheme (Lemma \ref{lemma-locally-closed-in-Noetherian}).
Hence $S$ is a sober, Noetherian topological space (use
Schemes, Lemma \ref{schemes-lemma-scheme-sober}).
Assume $S$ is not Jacobson to
get a contradiction. By
Topology, Lemma \ref{topology-lemma-non-jacobson-Noetherian-characterize}
there exists some non-closed point $\xi \in S$
such that $\{\xi\}$ is locally closed. This corresponds
to a prime $\mathfrak p \subset R$ such that (1) there exists
a prime $\mathfrak q$, $\mathfrak p \subset \mathfrak q \subset \mathfrak m$
with both inclusions strict, and (2) $\{\mathfrak p\}$ is open in
$\Spec(R/\mathfrak p)$. This is impossible by Algebra,
Lemma \ref{algebra-lemma-Noetherian-local-domain-dim-2-infinite-opens}.
\end{proof}







